\documentclass[11pt]{article}
\usepackage[margin=1in]{geometry}
\usepackage{amsmath,amssymb,amsthm,mathtools}
\usepackage{graphicx}
\usepackage{microtype}
\usepackage[hidelinks]{hyperref}

\newcommand{\R}{\mathbb R}
\newcommand{\norm}[1]{\left\lVert#1\right\rVert}
\newtheorem{theorem}{Theorem}
\newtheorem{proposition}{Proposition}

\title{No dimension-only stability constant\\for approximate
\(m\)-star-convexity}
\author{Candidate counterexample for arXiv:2207.08058, Problem 1}
\date{19 July 2026}

\begin{document}
\maketitle

\begin{center}
\fbox{\parbox{0.91\textwidth}{
\textbf{Status: candidate counterexample, likely valid.}
Problem 1 has a negative answer as literally stated: a stability constant
depending only on the dimension cannot be uniform in the parameter
\(m\in(0,1]\).  Any corrected constant must be at least
\(1/(1-m)\).}}
\end{center}

\section{Source and problem}

The source is Geanina Maria L\u{a}chescu and Ionel Roven\c{t}a,
\emph{The Hardy-Littlewood-P\'{o}lya inequality of majorization in the context
of \(\omega\)-\(m\)-star-convex functions}, arXiv:2207.08058 and
\emph{Aequationes Mathematicae} 97 (2023), 523--535 \cite{LR}.

For a convex set \(C\subseteq\R^N\), the source calls
\(\Phi:C\to\R\) \(\delta\)-\(\omega\)-\(m\)-star-convex when
\begin{equation}\label{eq:approx}
\begin{split}
 \Phi((1-\lambda)x+\lambda m y)
 \leq{}&(1-\lambda)\Phi(x)+m\lambda\Phi(y)\\
 &{}-m\lambda(1-\lambda)\omega(\norm{x-y})+\delta
\end{split}
\end{equation}
for all \(x,y\in C\) and \(\lambda\in(0,1)\).  The exact class is obtained
by deleting \(\delta\).

Problem 1 asks whether every such \(\Phi\) decomposes as
\(\Phi=\Psi+\Pi\), where \(\Psi\) is exactly
\(\omega\)-\(m\)-star-convex and
\begin{equation}\label{eq:askedbound}
 \norm{\Pi}_\infty\leq k_N\delta,
\end{equation}
with a positive constant \(k_N\) depending only on the dimension \(N\).

\begin{figure}[ht]
\centering
\includegraphics[width=\textwidth]{figures/open_problem_crop.png}
\caption{Definition 4 and Problem 1 on page 8 of the source paper.  The
source explicitly says that \(k_N\) depends only on \(N\).}
\label{fig:problem}
\end{figure}

\section{Proof intuition}

For \(m<1\), the coefficients on the right of the star-convexity inequality
sum to \(1-\lambda(1-m)<1\).  At the origin this forces every exact function
to be nonpositive.  The allowed additive error \(\delta\), however, permits a
positive constant of height \(\delta/(1-m)\).  Therefore the distance from
this approximate function to the exact class is at least
\(\delta/(1-m)\).  This diverges as \(m\) approaches one and rules out a
constant depending only on dimension.

\section{Quantitative obstruction and counterexample}

\begin{proposition}\label{prop:fixedm}
Fix \(N\geq1\), \(0<m<1\), and \(\delta>0\).  Among all stability constants
for Problem 1 at this fixed value of \(m\), any constant valid for every
convex domain and every zero-modulus approximately \(m\)-star-convex function
must be at least
\[
 \frac{1}{1-m}.
\]
\end{proposition}

\begin{proof}
Take \(C=\R^N\), \(\omega\equiv0\), and define
\begin{equation}\label{eq:constant}
 \Phi(x)=c:=\frac{\delta}{1-m}\qquad(x\in\R^N).
\end{equation}
For arbitrary \(x,y\in\R^N\) and \(\lambda\in(0,1)\), the right-hand side
of \eqref{eq:approx} is
\begin{align*}
 (1-\lambda)c+m\lambda c+\delta
 &=c-\lambda(1-m)c+\delta\\
 &=c+(1-\lambda)\delta\\
 &\geq c.
\end{align*}
This equals or exceeds
\(\Phi((1-\lambda)x+\lambda m y)=c\), so \(\Phi\) is
\(\delta\)-\(0\)-\(m\)-star-convex.

Now let \(\Psi:\R^N\to\R\) be exactly \(m\)-star-convex.  Substituting
\(x=y=0\) in its defining inequality gives
\[
 \Psi(0)\leq(1-\lambda+m\lambda)\Psi(0)
 =\bigl(1-\lambda(1-m)\bigr)\Psi(0).
\]
Since \(\lambda(1-m)>0\), it follows that \(\Psi(0)\leq0\).  Hence
\begin{equation}\label{eq:distance}
 \norm{\Phi-\Psi}_\infty
 \geq |\Phi(0)-\Psi(0)|
 \geq\frac{\delta}{1-m}.
\end{equation}
Thus no smaller universal factor can work at this \(m\).
\end{proof}

\begin{theorem}\label{thm:main}
There is no finite constant \(k_N\) depending only on \(N\) for which the
conclusion of Problem 1 holds for all \(m\in(0,1]\).  Consequently, Problem 1
as stated has a negative answer.
\end{theorem}

\begin{proof}
Fix \(N\) and suppose a finite dimension-only constant \(k_N\) exists.  Choose
\(m\in(0,1)\) so close to one that
\[
 \frac{1}{1-m}>k_N.
\]
The function \(\Phi\) from Proposition~\ref{prop:fixedm} is an admissible
approximate function, but \eqref{eq:distance} shows that every exact
\(m\)-star-convex \(\Psi\) satisfies
\[
 \norm{\Phi-\Psi}_\infty>k_N\delta.
\]
Therefore no decomposition satisfying \eqref{eq:askedbound} exists.
\end{proof}

\section{Verification, scope, and novelty}

The example uses the full-dimensional domain \(\R^N\) and the zero modulus.
It therefore avoids both the domain-closure ambiguity in the source definition
and any special choice of \(\omega\).  The two essential checks are the exact
calculation of the approximate defect and the origin substitution for exact
functions; both appear explicitly above.

The result does not exclude a stability theorem with a constant
\(k_{N,m}\) allowed to depend on \(m\).  Rather, Proposition~\ref{prop:fixedm}
shows that any such theorem covering the source's full class must satisfy
\[
 k_{N,m}\geq\frac{1}{1-m}.
\]
If the notation \(k_N\) was intended to suppress hidden dependence on a fixed
\(m\), the result is a sharp necessary-parameter obstruction instead of a
counterexample to that revised question.

Before promotion, the run indexes were searched using arXiv:2207.08058 and
the core terms for approximate \(m\)-star-convexity and Hyers--Ulam stability.
A bounded web/arXiv search used the exact Problem 1 sentence and close
stability phrases.  It found related results with extra hypotheses or
parameter-dependent bounds, but no answer to the dimension-only
\(\omega\)-\(m\) problem.  The source author's doctoral thesis,
\emph{Convex Analysis and Majorization Theory}, reproduces the problem in
Section 2.3.4 \cite{Thesis}.  This records the search bounds and is not an
exhaustive novelty guarantee.

Human review is recommended.  The main review issue is semantic rather than
algebraic: confirm that ``depends only on the dimension'' is intended to mean
uniformity in \(m\).

\begin{thebibliography}{9}
\bibitem{LR}
G. M. L\u{a}chescu and I. Roven\c{t}a,
\emph{The Hardy-Littlewood-P\'{o}lya inequality of majorization in the context
of \(\omega\)-\(m\)-star-convex functions},
Aequationes Math. 97 (2023), 523--535; arXiv:2207.08058.
\url{https://arxiv.org/abs/2207.08058}.

\bibitem{Thesis}
G. M. L\u{a}chescu,
\emph{Convex Analysis and Majorization Theory}, doctoral thesis, University
of Craiova, defended 2024 and confirmed 2025, Section 2.3.4.
\url{https://ndsas.ucv.ro/results/}.
\end{thebibliography}

\end{document}

